The campaign supports four mechanisms, each backed by matched comparisons and, where stochastic, by seed replicates.

\paragraph{1. Net section times ligament-size-dependent net-section strength.} The strength of a porous sheet is the minimum load-bearing solid fraction across the loading direction multiplied by a net-section strength that depends on the width and alignment of the load-bearing ligaments but not on pore shape: $\approx$30--32\,N/m (80\% of the intrinsic strength) for straight, load-parallel ligaments wider than $\sim$20\,\AA, 17--19\,N/m for round-pore ligaments of 10--13\,\AA, and 12--20\,N/m for random networks. Evidence: porosity independence of the parallel-slit strength, the linear net-section scaling of the square coarse mesh, the equality of square and round coarse pores at equal net section, the master plot of net-section strength versus ligament width, and the holdout predictions (Sec.~\ref{sec:holdout}). Mechanistically, the fine ligaments are weaker because a larger fraction of their bonds is perturbed by the two-coordinated pore edges (edge reconstruction and bond-order changes), an atomic-scale effect that has no continuum analogue at fixed geometry ratio.

\paragraph{2. Load-path orientation.} Ligaments or pores aligned with the load carry the stress without concentration; features across the load reduce the net section and concentrate stress at their ends. The same porosity yields 25.4 vs 4.0\,N/m for slits parallel vs perpendicular to the load, 19.0 vs 8.6\,N/m for ellipses, and 16.1 vs 12.0\,N/m for veins along vs across the load in nested meshes. This is the largest single lever in the design space and acts identically in single-scale and hierarchical structures.

\paragraph{3. Fibre-bundle load transfer sets the fracture mode.} After the first local failure, the released load is shared by the remaining parallel ligaments; failure is progressive only when the share transferred to each survivor stays below its strength margin, i.e. when there are many parallel paths of moderate strength dispersion. Few wide strips (three per cell) cascade even with 15\% dispersion (Stage~4), single flaws and clustered defects run in one event, whereas 8--10 rows of fine ligaments, jittered lattices, distributed vacancies and load-parallel-vein meshes fail row by row and retain 20--60\% of the peak load. Spatial localisation of the damage is a separate property: periodic single-avalanche failures are spatially delocalised, crack-like failures localised.

\paragraph{4. Lattice trapping of cracks.} Beyond $\approx$20\,\AA{} the strength of a cracked or holed sheet becomes independent of flaw size (18--19.5\,N/m for 20--50\,\AA{} flaws in 100--150\,\AA{} cells), because crack advance in the discrete lattice proceeds by trapped bond-breaking events with a barrier; stable bond-by-bond growth precedes the instability. Griffith scaling from the 10 and 20\,\AA{} cracks would have predicted 12--14\,N/m for the 40 and 50\,\AA{} cracks.

\paragraph{What hierarchy does.} A nested mesh behaves as the parallel combination of its finest level (which fails first) and its load-parallel coarse ligaments; at matched porosity it is never stronger than the coarse level alone, and never tougher than the fine level alone, unless the fine level is itself aligned with the load. Scale ratio, hierarchy depth and perpendicular veins are irrelevant or detrimental. Hierarchy is therefore not a mechanism in this model; alignment and ligament width are.
